\documentclass[12pt]{book}
\usepackage[a4paper,headheight=16pt,scale={0.8,0.85},hoffset=0.4cm]{geometry}

% Para escribir en castellano:
\usepackage[spanish]{babel}
\usepackage[utf8]{inputenc}

\usepackage[pdftex]{graphicx}
\usepackage{amsmath}
\usepackage{amssymb}
\usepackage{tipa}
\usepackage{textcomp}

\usepackage{fancyhdr}
\pagestyle{fancy}

%Comandos propios para simplificar escritura de capítulos
\usepackage{ejercicio}
\usepackage{nota}

\renewcommand{\chaptermark}[1]{\markboth{#1}{}}
\renewcommand{\sectionmark}[1]{\markright{\thesection\ #1}}
%Eliminando numeración y nombres de Capítulos y secciones
\addto\captionsspanish{
  \renewcommand\chaptername{}
}
\renewcommand{\thechapter}{}
\setcounter{secnumdepth}{-1} 

% Todas las imáenes está en el directorio img:
\newcommand{\imgdir}{img}
\graphicspath{{\imgdir/}}

\fancyhf{}
\fancyhead[LE,RO]{\bfseries\thepage}
\fancyhead[LO]{\bfseries\leftmark}
\fancyhead[RE]{\bfseries\leftmark}
\fancyfoot[R]{\includegraphics[height=1cm]{footer_logo}}
\renewcommand{\headrulewidth}{0.5pt}
\renewcommand{\footrulewidth}{0pt}
\addtolength{\headheight}{0.5pt}
\setlength{\footskip}{0.5in}
\renewcommand{\footruleskip}{0pt}
\fancypagestyle{plain}{%
 \fancyhead{}
 \renewcommand{\headrulewidth}{0pt}
}

%
%\parindent 0in
\parskip 0.05in

\usepackage{hyperref}
\hypersetup{
  pdftitle=Proyecto Nahual – Introducción a la programación,
  pdfauthor=Proyecto Nahual,
  pdfsubject=Libro de Introducción a la programación del Proyecto Nahual,
  pdfkeywords=programación nahual ruby,
  colorlinks,
  citecolor=black,
  filecolor=black,
  linkcolor=black,
  urlcolor=black,
  pdftex
}

%%%%%%%%%%%%%%%%%%%%%%%%%%%%%%%%%%%%%%%%%%%%%%%%%%%%%%%%%%%%%%%%%%%%%%%%%%%%%%%%
\usepackage{color}
\definecolor{gray75}{gray}{.75}

\definecolor{listinggray}{gray}{0.9}
\definecolor{lbcolor}{rgb}{0.9,0.9,0.9}
\usepackage{listings}
\usepackage{textcomp}
\lstset{
	backgroundcolor=\color{lbcolor},
	tabsize=4,
	rulecolor=,
	language=ruby,
        basicstyle=\scriptsize,
        upquote=true,
        aboveskip={1.5\baselineskip},
        columns=fixed,
        showstringspaces=false,
        extendedchars=true,
        breaklines=true,
        prebreak = \raisebox{0ex}[0ex][0ex]{\ensuremath{\hookleftarrow}},
        frame=single,
        showtabs=false,
        showspaces=false,
        showstringspaces=false,
        identifierstyle=\ttfamily,
        keywordstyle=\color[rgb]{0,0,1},
        commentstyle=\color[rgb]{0.133,0.545,0.133},
        stringstyle=\color[rgb]{0.627,0.126,0.941}
}
%Ejemplo de uso para incluir un archivo: \lstinputlisting{parameters.py}
%Ejemplo para codigo inline: \begin{lstlisting} \end{lstlisting}

\newlength{\PictHOffset}
\newlength{\PictVOffset}
\setlength{\PictHOffset}{1in}
\addtolength{\PictHOffset}{\hoffset}
\addtolength{\PictHOffset}{\oddsidemargin}

\setlength{\PictVOffset}{1in}
\addtolength{\PictVOffset}{\voffset}
\addtolength{\PictVOffset}{\topmargin}
\addtolength{\PictVOffset}{\headheight}
\addtolength{\PictVOffset}{\headsep}
\addtolength{\PictVOffset}{\topskip}
\addtolength{\PictVOffset}{-1\paperheight}
\begin{document}
\begin{titlepage}
   \noindent\hspace*{-\PictHOffset}%
   \raisebox{\PictVOffset}[0pt][0pt]{\makebox[0pt][l]{
      \includegraphics[width=\paperwidth,height=\paperheight]{portada}}
   }
   \vfill
     \begin{flushright}
     	v0.0.1
     	\end{flushright}
\end{titlepage}

\thispagestyle{empty}
\subsection*{\center \normalsize Alguna cosa sobre el libro}
 
\tableofcontents 
\mainmatter
\chapter{1er semana - ¿Qué es programar?}
\section{¿Qué es programar?}
Recuerdo claramente el momento en que escribí mi primer programa (y eso que yo tengo una memoria malísima), era un programa bastante complicado para alguien como yo. Pero como yo estudiaba matemática en la facultad se suponía que era bueno en cosas como ``pensar lógicamente''. Entonces fui a la sala de computadoras de la facu, armado solamente con un libro sobre programación y mi ego, me senté en una de las computadoras y empecé. Bueno, ``empecé'' no sería la palabra correcta. O tampoco ``programar". Digamos que me senté delante de la computadora y me sentí increíblemente tonto. Luego me sentí avergonzado. Y luego enojado hasta que finalmente me sentí pequeño: después de 8 horas desesperantes el programa estaba terminado y andando. El hecho de que funcionase no importaba mucho...era un momento de gloria.\\

Pero la pregunta que me quedo fue, ¿por qué fue tan difícil decirle a una computadora que haga algo medianamente complicado? Bueno, la parte difícil no es la de ``medianamente complicado'' sino la parte de ``decirle a una computadora``.\\

En cualquier comunicación entre humanos uno puede dejar muchas cosas sin decir, saltearse pasos y conceptos y entenderse. Claro que uno deja que la otra persona ``complete'' lo que falta. En realidad, uno tiene que saltearse cosas. Si no lo hiciésemos, ¡no podríamos hacer nada! El ejemplo típico es el  de armarse un sandwich de jamón, queso y mayonesa. Normalmente, uno diría tranquilamente: ``Che ¿me haces un sandwich de jamón y queso con mayonesa?'' En cambio, si se lo tenés que pedir a alguien que nunca hizo uno, como un robot, tendrías que decirle algo como:
\begin{enumerate}
  \item Agarrá 2 rodajas de pan (y dejá el resto)
  \item Sacá queso, jamón, mayonesa y un cuchillo
  \item Untá un poco de mayonesa en un pan
  \item Poné arriba de la mayonesa: 2 fetas de jamón y 1 de queso
  \item Poné la otra rodaja de pan arriba, poné todo en un plato y traémelo. ¡Gracias!
\end{enumerate}

Me imagino que esas instrucciones son suficientes para un niño de 7 años. Pero un niño de 7 años es bastante inteligente como para hacer eso. En cambio ¿qué habría que decirle a una computadora? Concentrémonos solamente en el primer paso:

\begin{enumerate}
  \item Agarrá 2 rodajas de pan (y dejá el resto)
  \begin{enumerate}
    \item Localizá la bolsa de pan lactal
    \item Andá a una mesa
    \item Dejá la bolsa de pan en la mesa
    \item Agarrá 2 rodajas de adentro de la bolsa
    \item \ldots
  \end{enumerate}
\end{enumerate}

Lamentablemente estas instrucciones no alcanzan para nada. Por ejemplo: ¿cómo es “localizá la bolsa de pan”? Tendríamos que tener una especie de base de datos en donde se asocien las cosas con su ubicación. La base de datos debería también tener las coordenadas del jamón, heladera, mesa, cuchillo, plato, yo mismo...
¡Ah! ¿Y qué pasaría si la bolsa de pan está en una panera? Entonces ¿habría que abrir la panera? O ¿por ahí está detrás de algo? ¿O arriba de la heladera? Entonces tendríamos solamente para localizar la bolsa de pan lactal:
\begin{enumerate}
  \item Buscá las coordenadas de la panera
  \item Si la bolsa de pan está en la panera
  \begin{enumerate}
    \item Abrí la panera
    \item Meté la mano dentro de la panera
    \item Agarrá la bolsa de pan lactal
    \item Sacá la mano de la panera
    \item Cerrá la panera
  \end{enumerate}
  \item Si la bolsa de pan está arriba de la heladera
  \begin{enumerate}
    \item Andá hasta la puerta de la heladera
    \item Estirá la mano hacia arriba de la heladera
    \item Agarrá la bolsa de pan lactal
    \item Sacá la mano de arriba de la heladera
  \end{enumerate}
\end{enumerate}

Y así sigue y sigue. Además ¿qué pasa si no hay cuchillo? ¿O si la mesa está llena de cosas y no hay lugar para preparar el sandwich? ¡Habría que rezar al Santo que te parezca más apropiado para que el pan esté en buen estado! Aún los pasos más sencillos, como abrir la panera, deberían ser explicados ¡y es por eso que todavía no hay robots que hagan sandwich para nosotros! No es que los robots no puedan ser construidos ¡sino que no podemos programarlos para que nos hagan un sandwich! Y eso es porque, aunque no lo parezca, preparar un sandwich de jamón, queso y mayonesa es algo muy difícil de describir (pero fácil de hacer para criaturas inteligentes como nosotros, los humanos) y las computadoras sólo son buenas para las cosas que son fáciles de describir. Y por eso para mí fue tan difícil hacer mi primer programa: ¡las computadoras son más tontas de lo que creía!

\section{Hacer un sandwich es programar}
Cuando uno le explica a alguien como hacer un sandwich, el trabajo es mucho más fácil porque esa persona ya sabe lo que es un sandwich. Es común que la gente sepa de “sandwichería” y por eso esa persona puede llenar los espacios de nuestra pobre explicación. El paso 3 dice “Untá un poco de mayonesa en un pan”, pero no dice que no lo unte en los dos lados. Tampoco dice cuánta mayonesa ni de qué manera untarla (si poner más en el medio, o en los costados) ni aclara que debe usar el cuchillo para hacerlo y no las manos. Uno asume que esas cosas ya se saben. Creo que esta realidad ayudaría a explicar qué quiere decir programar o por lo menos dar una definición informal.\\

Programar es decirle a tu computadora cómo hacer algo. Las tareas grandes tienen que ser divididas o descompuestas en tareas más chicas, las cuales tienen que ser descompuestas en tareas aún más chicas y así seguir hasta que las tareas sean tan chiquititas y simples que no tienen que ser divididas porque la computadora sabe hacerlas.\\

Estas tareas son cosas muy básicas como imprimir un texto en la pantalla o sumar 2 números. El problema más grande que tuve cuando estaba aprendiendo a programar fue que yo estaba aprendiendo al revés. Yo ya sabía lo que quería que la computadora hiciese y trataba de ir dividiendo eso hasta llegar a tareas que de la computadora entendiese. Muy mala idea. En realidad yo no sabía lo que realmente la computadora podía hacer, entonces no sabía bien en que tareítas descomponer mi tarea grande o lo que la computadora tenía que hacer.\\
 	
Por eso en Nahual, vas a aprender a programar en forma diferente. Primero vas a aprender estas tareítas básicas que la computadora puede hacer (algunas nomás) y después vas a encontrar tareas simples que puedan ser descompuestas en esas tareítas. Tu primer programa va a ser tan fácil que solamente te llevara menos de un minuto.

\section{Lenguajes de Programación}
Para poder decirle algo a la computadora tenemos que usar un lenguaje diferente. Un lenguaje de programación es parecido al lenguaje que usamos para hablar, en el sentido que está formado por cosas básicas como palabras (verbos y sustantivos) y por maneras de combinar esas palabras para dar sentido (oraciones, párrafos, novelas). Hay muchos lenguajes para elegir (C, Java, Ruby, Perl, \ldots) y cada uno tiene su conjunto de elementos básicos. Algunos tienen más que otros. Ruby tiene un conjunto sencillo pero elegante y es uno de los más fáciles para aprender así que vamos a usar ese.\\
 
Probablemente la mejor razón para usar Ruby es que los programas en Ruby suelen ser cortos. Por ejemplo aquí hay un pequeño programa en Java.

\begin{lstlisting}[language=java]
public class Hola Mundo {
  public static void main (String [] args) {
    System.out.println("Hola Mundo");
  }
}
\end{lstlisting}

Y acá esta el mismo programa pero en el lenguaje Ruby:

\begin{lstlisting}
  puts 'Hola Mundo'
\end{lstlisting}

Este programa, en ambos casos, escribe las palabras ``Hola Mundo'' en la pantalla de la computadora. En la versión en Ruby uno podría adivinar, pero en la versión de java adivinar esto no es tan sencillo.

\section{El arte de programar}
Una parte importante de programar es, claramente, que el programa haga lo que se supone que tiene que hacer y no otra cosa. En otras palabras, el programa no debería tener errores o \emph{bugs} (se pronuncia bags). Claro que, cuando uno se concentra mucho en no tener errores, se pierde la gran mayoría de las cosas buenas de la programación.\\

Programar no es solamente acerca del producto final o del programa sino acerca del proceso que lleva a él. Los programas, a diferencia de cosas más concretas como una casa, no se construyen en un orden tan rígido sino que los programadores charlan sobre él, piensan, discuten, escriben un poco, lo arreglan, lo estiran, juegan con él, lo cambian, lo mejoran, lo prueban, lo borran, lo vuelven a escribir, \ldots

\section{Un programa no se arma: crece}
Dado que un programa siempre está creciendo y evolucionando, siempre está cambiando y tiene que estar escrito pensando en que va a cambiar. Sé que esto no puede quedar muy claro hablando en términos prácticos pero este tema ira surgiendo durante todo el libro. Probablemente, la primera regla de un buen programador es evitar duplicar código, esto es ``No vuelva a escribir lo mismo". Si la vemos de otra manera, es como la ley del mínimo esfuerzo. Un buen programador tiene que ser un poco vago, pero no ser vago de no hacer nada, sino ser vago en forma inteligente ¡para no tener que hacer nada después! ¡Ahórrate trabajo siempre que puedas! Si hacer cambios ahora significa ahorrarte trabajo más adelante ¡hacelo!\\

Trata de que tu programa haga lo menos posible para hacer lo que tiene que hacer. De esta manera programar no solo es mucho más interesante (porque es muy aburrido hacer lo mismo una y otra vez) sino que, además ¡el programa tiene menos errores y es más rápido! ¡Es una situación donde todos ganan!\\

De cualquier manera, la idea es la misma: tratá de que tus programas sean flexibles. Así, cuando haya que hacer los cambios (y te aseguro que siempre vienen) entonces vas a pasar menos y mejor tiempo haciéndolos.\\

\section{Ejercicios}
\ejercicio{Mostrar un programa}{
Les mostramos un programa hecho en Ruby de los robots.
En este capítulo proponemos un role-play con un robot. Los chicos tienen que escribir los pasos del robot y el robot (uno de los instructores) tiene que repetirlo.}
\chapter{2da semana - Números y letras}

Aca citamos alguna pagina \cite{algo} y algun libro \cite{otro} :D
\chapter{3er semana - Variables y Compoteras}
\chapter{4ta semana - ¿Qué pasa si...?}
\chapter{5ta semana - Programar Jugando / La Mansión Maníaca}
\chapter{6ta semana - Nuestro primer juego (de texto)}
\nota{Introducción}{Ya somos capaces de controlar el flujo de un programa, o al menos estamos empezando a hacerlo, además ya entendemos lo que son las instrucciones, las condiciones, es hora de empezar a crear nuestros propios juegos!}
\nota{Objetivo}{El objetivo de esta semana es realizar nuestro primer juego, fortaleciendo nuestro conocimiento sobre cómo interactuar con la computadora}
\nota{Cómo lo vamos a hacer}{Vamos a desarrollar el juego utilizando todo lo que aprendimos hasta ahora interactuando con el usuario mediante las sentencias \emph{puts} y \emph{gets}, y luego manejando el flujo del programa con la sentencia \emph{if}}

\section{Nuestro primer juego}
La idea de nuestro primer juego es hacer un juego de consola, como el que jugamos el otro día. Sólo que el nuestro va a ser sólo en modo texto ¿qué significa modo texto? la aburrida pantallita negra que venimos viendo todos los días... si, si, es tristísimo ¡pero no tanto, che! Ya le vamos a poner dibujitos, no se preocupen.\\

Entonces, vamos a plantear un problema y las opciones que tiene nuestro jugador. En base a lo que él elija, irán sucediendo diferentes situaciones que, básicamente, nos llevarán a una nueva situación.

\section{Introducción}
El joven Charly se despierta una mañana y nota que es muuuy tarde y que hoy, exactamente hoy, es la prueba final de matemática con la profesora Marta, famosa por desaprobar a todos (y sobre todo a los que llegan tarde).\\

Charly salta de la cama y se da cuenta de que tiene casi toda la ropa puesta, con lo cual no necesita vestirse, pero... ¡no encuentra las Zapatillas!

\subsubsection{Situación 1}
Charly necesita conseguirse el par de zapatillas para poder salir corriendo. Entonces, en esta secuencia la idea es que se consiga las zapatillas.\\

Ahora deberíamos plantear las opciones que tenga Charly, se me ocurren por ejemplo:
\begin{itemize}
  \item Ver abajo de la cama
  \item Buscar entre la ropa tirada en la pieza
  \item Fijarse dentro del armario
  \item Seguir durmiendo
\end{itemize}

Entonces ¿qué hacemos con las respuestas? digamos que se elige la opción “Ver abajo de la cama”. La Compu le debería responder “No hay nada abajo de la cama”... y volver a presentar la introducción.\\

Si elije la siguiente opción (“Buscar entre la ropa tirada en la pieza”), debería decirle que tiene una zapatilla ¿qué desea hacer?:
\begin{itemize}
  \item Agarrarla.
  \item Dejarla tirada.
\end{itemize}

Cualquiera de las dos opciones debería llevarlo de nuevo a la introducción... Sólo que, si decidió agarrar la zapatilla (como me imagino que hará cualquier ser humano normal), entonces deberíamos guardar en algún lugar la zapatilla agarrada.\\

Si elige “Fijarse dentro del armario” debería suceder lo mismo que en el caso anterior, ya que, como se estarán imaginando, dentro del armario está la otra zapatilla que busca Charly para poder ir a rendir matemática. ¿Qué hacemos entonces? mostramos de nuevo lo mismo:
\begin{itemize}
  \item Agarrarla.
  \item Dejarla en el armario.
\end{itemize}

Si el pibe es bastante piola y como corresponde ya tiene la otra zapatilla, entonces se deberá mostrar el siguiente mensaje:
\begin{console-output}
FIN DEL JUEGO: CHARLY SALIO CORRIENDO PARA EL COLEGIO CON SUS DOS ZAPATILLAS!
\end{console-output}

Por supuesto que si elige la última opción, o sea,  Seguir durmiendo, entonces deberá decirle
\begin{console-output}
FIN DEL JUEGO: MARTA MANDO A CHARLY A RENDIR EN DICIEMBRE!
\end{console-output}
\chapter{7ma semana - ¿Qué aprendimos hasta ahora?}
\chapter{8va semana - Tratemos de no escribir tanto}
\nota{Introducción}{Después de repasar todo lo que fuimos aprendiendo y del examen, ahora podemos empezar a aprender algunas herramientas que hacen más sencillo programar}
\nota{Objetivo}{El objetivo de esta semana es aprender sobre los bucles o \emph{loops} que nos permiten ejecutar sentencias que se repiten}
\nota{Cómo lo vamos a hacer}{Vamos a ir mostrando con ejemplos y ejercicios simples las diferentes formas de repetir que tenemos en Ruby}

\section{Control de Flujo}
Ahora vamos a manejar un poco el flujo del programa con cosas más divertidas que solo preguntar cosas de a una. Las palabras claves de esta clase son:

\subsection{Iterando}
Uh ¡se puso difícil la cosa! Esto de iterar suena raro pero van a ver que es una pavada... y no se asusten, porque encima les tengo que contar como se dice en ingles: \emph{loop}. Raro ¿no? Bueno, no tanto. Les cuento un poco cuándo se usa.\\

Resulta que cuando tenemos que repetir muchas veces una acción, como mostrar los números de uno a 20, es medio aburrido hacer:

\begin{center}
\begin{tabular}{|c|c|}
\hline
\rowcolor[gray]{0.9}Instrucciones a la Computadora & Lenguaje Humano \\
\hline
\begin{lstlisting}[language=ruby]
puts 1
puts 2
puts 3
puts 4
puts 5
puts 6
...
...
puts 19
puts 20
\end{lstlisting} & \begin{lstlisting}
mostra el uno
mostra el dos
mostra el tres
mostra el cuatro
mostra el cinco
mostra el seis
...
...
mostra el 19
mostra el 20
\end{lstlisting} \\
\hline
\end{tabular}
\end{center}

Y así...hasta mi abuela que es re pilas se aburre y se va de viaje a Mar del Plata con los jubilados. Entonces, ¿cómo podemos hacer con la compu? Podemos decirle algo así:

\begin{center}
\begin{tabular}{|c|c|}
\hline
\rowcolor[gray]{0.9}Instrucciones a la Computadora & Lenguaje Humano \\
\hline
\begin{lstlisting}[language=ruby]
numero = 0;
while numero < 20
  puts numero
  numero = numero + 1
end
\end{lstlisting} & \begin{lstlisting}
define la variable numero en cero
mientras numero es menor a veinte
muestra el valor de numero
suma en uno a numero
repetir eso
\end{lstlisting} \\
\hline
\end{tabular}
\end{center}

No sé ustedes, pero mi abuela y yo preferimos hacerlo así porque, además, si en vez de veinte son cien, es más difícil y para algo tenemos la compu ¿no?

\subsection{La estructura estándar de un loop}
Analicemos ahora un poco más en detalle el ejemplo anterior, veamos la primera línea de código:

\begin{lstlisting}
numero = 0
\end{lstlisting}
Esto, que parece tan trivial, realmente tiene mucha importancia: estamos inicializando la variable \emph{numero} en cero. Hacemos esto porque queremos contar desde 0 hasta 19, ahora ¿qué pasaría si quisiéramos contar desde 1? Simplemente cambiaríamos el valor de número por un valor inicial en uno:

\begin{lstlisting}
numero = 1
\end{lstlisting}

Veamos ahora la siguiente línea
\begin{lstlisting}
while numero < 20
\end{lstlisting}

En esta línea tenemos algo que es clave, es la condición que se tiene que cumplir para que la compu siga iterando, en nuestro caso el número que se muestra llegaría hasta el 19, pero ¿Qué pasaría si queremos que llegue hasta el 20? ¿qué tenemos que cambiar? La respuesta es: la condición, tenemos que hacer una que incluya el 20, ya que menor que 20 no lo incluye, pero menor o igual a 20 si.
\begin{lstlisting}
while numero <= 20
\end{lstlisting}

Acá estaríamos diciéndole hace mientras el numero sea menor o IGUAL a 20, y esa condición si la cumple el 20.\\

La siguiente línea es bastante simple, es imprimir el número.
\begin{lstlisting}
puts numero
\end{lstlisting}

Esta era fácil, pero la que se viene es otra importantísima: es el incremento de la variable \emph{numero} ¿Qué significa eso? significa el pequeño paso que voy a avanzar en cada uno de las iteraciones. En este caso, como estoy sumando los números consecutivos de a uno:
\begin{lstlisting}
numero = numero + 1
\end{lstlisting}

\subsection{Iterando con preguntas}
Bueno, pero ahora tenemos que repetir cosas con preguntas y palabras, ¿Qué se pensaban? ¿Qué era joda esto? Vamos a ver que no es tan difícil ¡al contrario! les va a resultar tan fácil que se van a terminar yendo de vacaciones con la abuela.\\

Miremos el ejemplo de despedida:
\begin{lstlisting}
input = ' '
while input != 'chau'
  puts input
  input = gets.chomp
end
puts 'Vuelvan prontoooo!!!'
\end{lstlisting}

Acá usamos bastantes cosas divertidas, como el “hacer mientras” que aprendimos recién y el \emph{gets} del primer programa interactivo. Se lo dejamos para que nos lo expliquen ustedes.

\section{Ejercicios}
\begin{ejercicio}{Variables}{
\begin{enumerate}
  \item Hacer un programa que muestre los números que van del 30 al 50
  \item Hacer un programa que muestre los números pares entre 1 y 70
  \item Hacer un programa que muestre los números impares entre 2 y 90
  \item Hacer un programa que muestre los números entre 1 y 70, pero al revés, o sea empezando por el 70, siguiendo por el 69, 68...
  \item Hacer un programa que ingresando un número lo eleve al cuadrado (ejemplo: 5 al cuadrado (2) = 25)
  \item Hacer un programa que ingresando un número lo eleve al cubo (ejemplo: 2 al cubo (3) = 8)
  \item Hacer un programa que ingresando un número lo eleve a cualquier potencia (pedir el número y la  potencia. Ejemplo: 2 a la cuarta = 16)
  \item Hacer un programa que escriba al revés una palabra dada (ejemplo: Mariano => onairaM)
  \item Hacer un programa que calcule el resto de una división (pedir el dividendo y el divisor. ejemplo: el resto de 7 dividido 2 es 1)
  \end{enumerate} 
}
\end{ejercicio}

\subsection{Ejercicio resuelto}
Solución del ejercicio de la palabra:
\lstinputlisting{code/cap08_ej_palabras.rb}

\chapter{9na semana - Repaso y Ejercicios}
\section{Guía de repaso}
\nota{Objetivo}{El objetivo de esta guía es ayudarte a buscar cómo resolver pedazos de ejercicios, trabajos prácticos o bien en cualquier momento que te quieras hacer un programa en Ruby, por eso el equipo de Nahual te armó este documento}
\nota{Como usar la guía}{La guía está escrita de una manera amigable para que te sea fácil buscar lo que querés, la mayoría de los títulos van a ser sentencias, y vamos a tratar de que tengan también ejemplos para que puedas usarlas}

La guía arranca de lo más fácil a lo más complejo, por lo tanto, si vos estás buscando algo como iteradores o condiciones que son las últimas cosas que vimos, es probable que las encuentres cerca del final, y si buscás cómo convertir datos o cómo transformar un numero en una letra y cosas así que vimos al principio, van a estar en la primer parte.

\subsection{Mostrar un mensaje en la pantalla}
Cuando un ejercicio nos dice, por ejemplo, que tenemos que poder elevar un numero a cualquier potencia, lo más básico es pensar que vamos a mostrarle al usuario un mensaje que diga, “ingrese un numero:”, esto se hace en Ruby con:
\begin{lstlisting}
puts 'Ingrese un numero'
\end{lstlisting}

\subsection{Mostrar un número en pantalla}
\begin{lstlisting}
puts 5
\end{lstlisting}

\subsection{Recibir un dato de la pantalla}
Cuando quiera guardar un dato que me ingresan, lo hago en una variable, por ejemplo numero, entonces tengo que hacer:
\begin{lstlisting}
puts 'Ingrese un numero'
numero = gets.chomp
\end{lstlisting}

Y ahí el valor de lo que el usuario escriba por pantalla queda en la variable numero, por ejemplo

\begin{lstlisting}
puts 'Ingrese su nombre'
nombre = gets.chomp
\end{lstlisting}
\subsection{Convertir un número a letra}
Muchas veces tenemos que escribir un mensaje, por ejemplo, que sea "el resultado es 2" y esto es bastante complicado porque 2 es un número, entonces hacemos lo siguiente

\begin{lstlisting}
2.to_s
\end{lstlisting}

\subsection{Convertir una letra a número}
Cuando recibimos los datos desde la pantalla son solo letras, entonces de golpe uno puede querer convertirlos a número, entonces hace:

\begin{lstlisting}
variable.to_i
\end{lstlisting}

\subsection{Hacer algo según una condición}
Si se cumple cierta condición tenemos que hacer algo, sino otra cosa:

\begin{lstlisting}
if (condicion)
  hacemos lo que se hace cuando la condicion es verdadera 
else
  hacemos lo que se hace cuando es falsa
end
\end{lstlisting}

\subsection{Repetir algo mientras se cumple una condición}
A veces queremos poder repetir algo mientras se cumple una cierta condición

\begin{lstlisting}
while{condicion}
 hacer lo que queremos hacer muchas veces
end
\end{lstlisting}

\section{Cuestionario chiquito}
\nota{Objetivo}{Necesitamos ir aprendiendo constantemente y para eso a veces es bueno poder repasar un poquito las cosas que aprendemos, como hoy es un día de repaso vamos a hacer un cuestionario cortito que nos sirve para saber cómo venimos, aclaramos que EL CUESTIONARIO NO LLEVA NOMBRE Y NO ES UNA PRUEBA solamente es para saber dónde estamos parados}

\subsection{Números y Letras}
\begin{enumerate}
  \item ¿Qué resultado devuelve hacer 2 + 3?
  \item ¿Qué resultado devuelve hacer 2.0 + 3?
  \item ¿Qué resultado devuelve hacer 2 + 3.5?
  \item ¿Qué resultado devuelve hacer 'Hola' + ' Mariano'?
\end{enumerate}

\subsection{Condiciones}
\begin{enumerate}
  \item ¿Qué resultado devuelve 3 == 3?
  \item ¿Qué resultado devuelve 3 > 3?
  \item ¿Qué resultado devuelve 3 >= 3?
  \item ¿Qué resultado devuelve 'Hola' >= 'Mariano'?
  \item ¿Qué resultado devuelve 'Hola' >= 'mariano'?
\end{enumerate}

\subsection{Funciones especiales}
\begin{enumerate}
  \item ¿Cómo transformamos de número a letra?
  \item ¿Cómo transformamos de mayúscula a minúscula?
  \item ¿Cómo transformamos de letra a número?
\end{enumerate}

\subsection{Variables}
\begin{enumerate}
  \item ¿Qué estamos haciendo cuando hacemos variable = 'Mariano'?
  \item ¿Qué devuelve el siguiente programita?
\end{enumerate}

\begin{lstlisting}
mi_nombre = 'Charly'
mi_apellido = 'Lizarralde'
puts 'Yo me llamo ' + mi_nombre + '  y mi apellido es ' + mi_apellido
\end{lstlisting}

\subsection{Loops}
En el siguiente programa:

\begin{lstlisting}
contador = 20
while contador < 30
  puts contador
  contador = contador + 2
end
\end{lstlisting}
\begin{enumerate}
  \item ¿Con qué valor se inicializa el contador?
  \item ¿Cuál es la condición?
  \item ¿En cuanto se incrementa el contador?
  \item ¿Qué devuelve el programa?
\end{enumerate}

\section{Ejercicios de repaso}
\nota{Objetivo}{El objetivo de esta guía es repasar algunos ejercicios que fuimos haciendo y juntarlos, sobre todo los últimos de condiciones y ciclos para que podamos seguir practicando ¡ya que ahora tenemos días para hacer ejercicios!}
\subsection{Loops}
\nota{Pistas}{Tenemos que definir las partes de un ciclo repetitivo o loop ¿se acuerdan? la condición, el valor inicial y como incrementamos la variable que vamos sumando}
\begin{enumerate}
  \item Hacer un programa que muestre los números que van del 30 al 50
  \item Hacer un programa que muestre los números pares entre 1 y 70
  \item Hacer un programa que muestre los números impares entre 2 y 90
  \item Hacer un programa que muestre los números entre 1 y 70, pero al revés, o sea empezando por el 70, siguiendo por el 69, 68.
\end{enumerate}

\subsection{Loops, pero con condiciones más difíciles}
\nota{Pistas}{Ahora vamos a tener que pensar un poco más cuándo terminamos de iterar y cómo calculamos}

Resulta que El Pichuleo, un joven atorrante de Entre Ríos, se fue a comprar ropa a la feria de la salada, pero solo llevó 100 pesos. Necesita comprar varias prendas y los precios que encuentra en la salada son:
\begin{itemize}
  \item Jean: 20 pesos
  \item Zapatillas Nike 40 pesos
  \item Calzoncillo de Boca: 5 pesos
  \item Remera de River: 50 pesos
  \item Remera de Independiente: 50 pesos
  \item Remera de Boca: 40 pesos
  \item Ojotas: 15 pesos
  \item Medias con la foto de Carrió: 13 pesos
  \item Medias de fútbol: 15 pesos
\end{itemize}

Entonces, El Pichuleo necesita un programita que le deje ir comprando y siempre le avise cuánta plata le queda. Ejemplo, por pantalla le aparece:

\begin{lstlisting}
Pichu, te quedan 100 pesos. Que te queres comprar?
1) Jean
2) Zapatillas Nike
3) Calzoncillo de Boca
4) Remera de River
5) Remera de Independiente
6) Remera de Boca
7) Ojotas 
8) Medias con la foto de Carrio
9) Medias de futbol
0) Salir
\end{lstlisting}
Si Pichu elige 2, nuestro programa hace lo siguiente:
\begin{lstlisting}
Pichu, te compraste unas zapatillas Nike
Pichu, te quedan 60 pesos. Que te queres comprar?
1) Jean
2) Zapatillas Nike
3) Calzoncillo de Boca
4) Remera de River
5) Remera de Independiente
6) Remera de Boca
7) Ojotas 
8) Medias con la foto de Carrio
9) Medias de futbol
0) Salir
\end{lstlisting}

Si se quiere comprar algo que no le alcanza el programa le dice:
\begin{lstlisting}
Pichu, no te alcanza la plata para eso
\end{lstlisting}
y sigue hasta que se elige 0 que es salir.
\chapter{10ma semana - Cajones llenos de cosas (listas)}
\chapter{11va semana - Funciones}
\chapter{12va semana - Jugando con las funciones de otro}
\chapter{13va semana - Fin}
En esta última semana, les presentamos Ruby On Rails (para hacer páginas web con Ruby), les damos un CD con el instalador de Ruby, ejemplos y juegos... y ¡un diploma por haber terminado el curso! 
\end{document}